\section{결론}
\label{sec:conclusion}

기본 CoC는 주행 장면에서 선택한 행동의 이유와 목표를 설명하지만, 같은 목표를 달성하는 여러 궤적 중 어느 범위까지 허용되는지는 충분히 나타내지 못할 수 있다. 본 논문은 이러한 허용 조건을 보완하기 위해 기본 CoC와 실행 가드를 하나의 문맥으로 결합한 \sccoc{}를 제안하였다. 또한 같은 장면, 행동 목표 및 궤적 후보를 유지한 채 실행 가드만 바꾸는 쌍대 계약 평가를 구성하여, 모델이 장면 유형이나 행동 목표가 아니라 실제 가드의 차이를 읽고 궤적을 선택하는지 확인하였다.

Qwen3-VL-2B-Instruct를 미세 조정한 표준 시간 평가에서 제약 정보가 없는 기본 CoC의 가드 위반률은 $0.319\pm0.052$였으나, 제약 통합 CoC에서는 $0.000\pm0.000$이었다. 이때 가드 준수 목표 완료율과 계약 쌍 정확도는 모두 $1.000$이어서 학습과 같은 문장 틀ㆍ수치 범위에서는 제약에 따른 후보 전환을 학습할 수 있었다. 그러나 학습에 없던 대기시간을 사용한 미관측 시간값 평가에서는 후보 정확도 0.635, 가드 위반률 0.351, 가드 준수 목표 완료율 0.649, 계약 쌍 정확도 0.306으로 크게 낮아졌다. 저장 어댑터의 재평가는 모든 원본 지표를 정확히 재현하였다. 표준 시간 과제의 후보 정확도는 단색 영상에서 0.670, 상충 구역이 비기 시작한 시각이 실제와 다른 충돌 영상에서 0.608로 낮아져 합성 연속 장면 영상의 시간 정보에 대한 의존성을 보였지만, 다중 주행의 표준ㆍ미관측 수치 평가는 영상 개입에도 정확도 1.000을 유지하여 문장 정보만으로 풀 수 있었다. 정지ㆍ차선 변경의 복합 표현 변화 평가에서는 위반률 $0.010\pm0.010$, 계약 쌍 정확도 $0.896\pm0.127$이었고 관찰된 오답은 m--cm 표현과 함께 나타났다. \revfinal{이 평가는 요인별 실험이 아니므로 단위 변환을 포함한 특정 요인을 성능 저하의 주된 원인으로 확정할 수 없다.} 따라서 결과는 합성 표준 분포에서의 제약 기반 후보 선택과 과제별 영상 사용을 지지하지만, 새로운 수치 관계의 일반적 해석을 지지하지는 않는다.

\revtwo{제안 방법의 성능은 모델 외부 검증 결과와 상호 보완 지표를 함께 볼 때 가장 분명하다. 표준 정적ㆍ시간ㆍ정지/차선 변경 평가에서 제약 통합 CoC의 가드 위반률은 모두 0.000이고 계약 쌍 정확도는 모두 1.000이었다. 정적 과제의 목표 완료율과 시간ㆍ다중 주행 과제의 가드 준수 목표 완료율도 1.000이므로, 이 결과는 항상 정지하거나 목표를 포기해 얻은 낮은 위반률이 아니다. 또한 쌍대 계약의 두 가드를 모두 맞혔으므로 같은 장면 답을 반복한 결과도 아니다. 외부 검증기는 모델의 자기 설명이 아니라 저장된 정지 위치, 속도, 감속도, 가가속도, 간격 및 진입 시각을 계약 경계와 비교해 이 판정을 재현 가능하게 계산한다. 다만 미관측 시간값과 복합 표현 변화에서 성능이 낮아졌으므로 이러한 우수성은 표준 합성 평가에 한정된다.}

본 연구의 결과는 합성 장면과 미리 정한 궤적 후보를 사용한 작은 VLM의 선택 실험에 한정된다. 영상 절제는 정보 기여를 과제별로 분리했지만 글자가 그려진 합성 영상만 사용했고, 정적 실험은 한 난수 초기값만 사용했으며, 세 난수 초기값 반복은 데이터 생성과 학습 난수를 함께 바꿨다. 검증 규칙도 합성 자료 생성 규칙과 같은 과제 정의에 속한다. 따라서 실제 도로의 시각적 근거화, Alpamayo-R1의 연속 궤적 또는 실제 차량의 안전성을 직접 입증하지 않는다. 또한 본 논문은 이미 주어진 실행 가드의 학습과 준수를 다루었을 뿐, 실제 CoC와 주행 환경으로부터 필요한 제약 조건, 수치 경계 및 해제 조건을 생성하는 방법은 제시하지 않았다. 별도 소형 신경망의 보조 폐루프 실험에서도 동적 계약과 차단장치를 함께 사용하면 위반이 줄었지만, 계약 갱신이 늦어지면 다시 증가하였다. 이는 실행 검사의 역할을 예시할 뿐 Qwen 결과의 폐루프 재현이나 안전 보장은 아니다.

\revone{폐루프 보조 분석에서 차단장치는 해제 후 재위험의 위반률과 충돌률을 각각 0.166과 0.125 낮췄지만 episode당 9.935회의 개입, 진행거리 0.030 m 감소 및 평균 절대 가가속도 0.306 증가가 동반되었다. 이미 이진 안전지표가 0인 phase-complete bank에서는 episode당 1.269--2.539회 개입하고도 추가 안전 이득이 없었다. 따라서 실제 적용 가능성은 안전지표와 intervention cost를 함께 평가해야 한다.}

향후 연구에서는 먼저 자연 영상과 정보량을 맞춘 시각 대조군, 그리고 수치 요인의 개별 변형을 사용하여 시각 정보, 문장 표현과 수치 비교의 기여를 더 세밀하게 분리해야 한다. 그다음 CoC의 객체ㆍ원인ㆍ요구 행동ㆍ시간 관계를 교통 법규, ODD, 시스템 안전 요구사항, 차량 동역학 및 센서 불확실성의 근거와 연결하여 제약 후보를 생성할 필요가 있다. 생성된 후보는 값과 단위, 유지ㆍ해제 조건, 적용 범위, 출처 및 불확실할 때의 대체 행동을 포함하는 구조화 계약으로 표현하고, 전문가 검토와 반례 기반 시뮬레이션을 통해 반복적으로 보정해야 한다. 이후 검증된 계약을 대형 자율주행 VLM의 연속 궤적에 적용하고, 최신 센서 상태를 이용하는 모델 외부 궤적 검증기 및 차단장치와 결합하여 계약의 충분성, 과도한 보수성 및 실제 환경에서의 적용 가능성을 평가할 필요가 있다.
